\documentclass[12pt,a4paper]{article}

% ── Paquets ───────────────────────────────────────────────
\usepackage[utf8]{inputenc}
\usepackage[T1]{fontenc}
\usepackage[french]{babel}
\usepackage{geometry}
\geometry{margin=2.5cm}
\usepackage{graphicx}
\usepackage{float}
\usepackage{enumitem}
\usepackage{textcomp}
\usepackage{hyperref}
\usepackage{array}
\usepackage{booktabs}
\usepackage{longtable}
\usepackage{xcolor}
\usepackage{listings}
\usepackage{fancyhdr}
\usepackage{titlesec}
\usepackage{multirow}
\usepackage{tabularx}
\usepackage{makecell}
\usepackage{colortbl}  % pour \rowcolor

% ── Liens ─────────────────────────────────────────────────
\hypersetup{
    colorlinks=true,
    linkcolor=blue,
    urlcolor=cyan,
    pdftitle={Rapport TFE --- SENTYS},
}

% ── Code ──────────────────────────────────────────────────
\lstset{
    basicstyle=\ttfamily\small,
    breaklines=true,
    frame=single,
    backgroundcolor=\color{gray!10},
    keywordstyle=\color{blue},
    commentstyle=\color{gray},
}

% ── En-tête / pied de page ────────────────────────────────
\pagestyle{fancy}
\fancyhf{}
\rhead{HE202204 EPHEC}
\lhead{Sentys}
\cfoot{\thepage}
\renewcommand{\headrulewidth}{0.4pt}

% ── Interligne ────────────────────────────────────────────
\renewcommand{\baselinestretch}{1.15}

% ==========================================================
\begin{document}

% ── Page de garde ─────────────────────────────────────────
\thispagestyle{empty}
\begin{center}

\includegraphics[width=0.75\textwidth]{LOGO._EPHEC.png}

\vspace{2cm}
{\LARGE \textbf{Rapport de Travail de Fin d'Études}} \\[0.6cm]
{\large \textbf{SENTYS}} \\[0.3cm]
{\large Système de Surveillance Résidentiel Résilient} \\[0.3cm]

\vspace{2cm}

{\large
\textbf{Théo Mertens} \\[0.4cm]
}

\vspace{1.5cm}

{\large
% CORRECTION : "Bachelier en Technologies de l'Informatique" et non "en Informatique"
\textbf{Formation :} Bachelier en Technologies de l'Informatique --- BAC 3 \\[0.2cm]
\textbf{Institution :} EPHEC Louvain-la-Neuve \\[0.2cm]
\textbf{Année académique :} 2025 -- 2026 \\[0.2cm]
}

\end{center}

\newpage
\thispagestyle{empty}

% ── Table des matières ────────────────────────────────────
\tableofcontents

\newpage

% ==========================================================
\section{Remerciements}
% ==========================================================

Je tiens en premier lieu à remercier ma famille pour son soutien constant et ses précieux conseils tout au long de ces années d'études.

Mes remerciements vont également à mes amis de classe pour leurs échanges enrichissants, leur solidarité et les nombreux moments de partage qui ont jalonné ce parcours.

Je remercie ensuite Monsieur Dewulf, mon rapporteur de TFE, pour son suivi attentif,
sa disponibilité et les orientations et remarques qu'il m'a apportées lors de l'élaboration de ce travail.

Je souhaite enfin exprimer ma gratitude à l'ensemble du corps professoral de l'EPHEC Louvain-la-Neuve, dont l'enseignement m'a fourni les fondements techniques et méthodologiques nécessaires à la réalisation de ce projet.


\newpage

% ==========================================================
\section{Introduction}
% ==========================================================
% CORRECTION : section 2 disposait d'un unique 2.1 sans 2.2 ni 2.3.
% Trois sous-sections sont maintenant justifiées.

\subsection{Description du projet}

SENTYS est un système de surveillance résidentiel complet, entièrement local et autonome, développé dans le cadre du TFE de Bachelier en Technologies de l'Informatique à l'EPHEC.

Le projet répond à un besoin concret : sécuriser une habitation en zone semi-isolée, disposant d'une connexion Internet instable, avec des exigences fortes en matière de confidentialité des données et d'autonomie opérationnelle. Le système doit fonctionner sans dépendance à aucun service cloud externe, et doit continuer à opérer en cas de coupure WiFi ou de panne électrique.

Le nom SENTYS fait référence à une sentinelle et évoque un système permettant de surveiller son domicile.

\subsection{Contexte et enjeux de souveraineté numérique}
% CORRECTION : la souveraineté numérique était mentionnée au 3.1.2 seulement.
% Elle constitue pourtant un enjeu central qui doit être posé dès l'introduction.

Le marché de la vidéosurveillance résidentielle est aujourd'hui largement dominé par des plateformes cloud commerciales (Ring, Nest, Arlo). Ces solutions imposent, par nature, le transit et le stockage des flux vidéo sur des serveurs tiers, souvent localisés hors de l'Union Européenne. Cette dépendance soulève trois problèmes fondamentaux :

\begin{enumerate}
    \item \textbf{Souveraineté des données :} les images captées dans un domicile privé sont confiées à des entreprises dont les conditions d'utilisation peuvent évoluer, être piratées, ou être soumises à des injonctions légales étrangères.
    \item \textbf{Dépendance opérationnelle :} en l'absence de connexion Internet, ces systèmes deviennent totalement inopérants, ce qui est particulièrement problématique en zone à connectivité instable.
    \item \textbf{Coût récurrent :} l'accès aux fonctionnalités essentielles (archivage, détection de mouvement) est généralement conditionné à un abonnement mensuel, représentant une charge financière permanente.
    \item \textbf{Accessibilité en cas de coupure secteur :} une coupure électrique met hors ligne le routeur Internet du domicile, rendant tout accès distant impossible même si le système de surveillance dispose d'une alimentation de secours. Pour garantir la continuité de l'accès à distance, une connexion 4G de secours est nécessaire, représentant un coût matériel et d'abonnement supplémentaire.
\end{enumerate}


SENTYS s'inscrit dans une démarche alternative et citoyenne : concevoir un système où l'intégralité des traitements, des données et du stockage reste sous le contrôle exclusif de l'utilisateur, sans aucun intermédiaire cloud.

\subsection{Structure du rapport}

Ce rapport est structuré en neuf sections. Après la présentation du cahier des charges (section~3) et l'analyse business (section~4), l'analyse technique (section~5) détaille les choix architecturaux et technologiques retenus parmi les solutions candidates. La stratégie de validation (section~6) décrit les tests effectués, y compris le déploiement en conditions réelles chez le client. La section~7 présente la mise en œuvre concrète, les fonctionnalités développées et les défis rencontrés. Les aspects législatifs (section~8) traitent des obligations légales liées à l'installation de caméras en Belgique. La conclusion (section~9) dresse le bilan du projet et ses perspectives d'évolution.

% ==========================================================
\section{Méthodologie et Planification}
% ==========================================================

\subsection{Approche de travail}

Le projet a été conduit selon une approche itérative, organisée en grandes phases successives. Chaque phase a donné lieu à un livrable intermédiaire discuté avec le promoteur, permettant de valider les choix effectués avant de passer à l'étape suivante. Cette méthode a permis de détecter tôt les écarts entre les besoins exprimés et la solution en cours de construction, notamment lors de la transition matérielle vers une architecture distribuée de nœuds sous Linux.

\subsection{Calendrier de travail}

\begin{table}[H]
\centering
\small
\renewcommand{\arraystretch}{1.3}
\begin{tabular}{|l|l|p{7cm}|}
\hline
\textbf{Période} & \textbf{Phase} & \textbf{Activités principales} \\
\hline
Septembre -- Octobre 2025 & Cadrage & Définition des besoins avec le client, rédaction du cahier des charges, premiers choix technologiques \\
\hline
Octobre -- Novembre 2025 & Analyse & Architecture système, modélisation de la base de données, comparaison des solutions matérielles et logicielles \\
\hline
Novembre -- Décembre 2025 & Développement (Phase 1) & Mise en place du serveur central, authentification JWT, pipeline de streaming vidéo, dashboard \\
\hline
Janvier 2026 & Développement (Phase 2) & Agent nœud caméra, mode hors ligne, synchronisation des clips, notifications push \\
\hline
Février 2026 & Développement (Phase 3) & Boîtier de contrôle physique, PWA, CI/CD, accès distant sécurisé \\
\hline
Mars 2026 & Tests et validation & Tests de résilience (coupures WiFi/secteur), tests de sécurité, déploiement chez le client \\
\hline
Avril -- Mai 2026 & Rédaction & Finalisation du rapport TFE, corrections suite aux relectures \\
\hline
\end{tabular}
\caption{Calendrier de réalisation du projet SENTYS}
\end{table}

\subsection{Échanges avec le client}

Le suivi du projet a été assuré par des échanges réguliers avec le client tout au long de l'année académique :

\begin{itemize}[label=\textbullet]
    \item \textbf{Réunion de cadrage (février 2025) :} première rencontre formelle pour identifier les besoins, les contraintes (zone semi-isolée, instabilité réseau, refus du cloud) et définir le périmètre du projet.
    \item \textbf{Points d'avancement mensuels :} présentation des développements en cours, recueil du feedback sur l'interface et les fonctionnalités, ajustement des priorités.
    \item \textbf{Validation du dashboard (Mi Avril 2026) :} démonstration en direct de l'interface web et du flux vidéo en temps réel. Le client a validé l'ergonomie générale et demandé l'ajout du boîtier de contrôle physique comme point d'entrée principal.
    \item \textbf{Déploiement de test (Fin avril 2026) :} installation du système complet sur site pendant plusieurs jours pour valider le comportement en conditions réelles (coupures réseau, autonomie batterie, qualité vidéo). Les résultats sont détaillés en section~\ref{sec:validation-terrain}.
\end{itemize}

% ==========================================================
\section{Cahier des Charges}
% ==========================================================

Cette section détaille les spécifications fonctionnelles et les contraintes opérationnelles du projet, visant à répondre aux besoins de souveraineté numérique et de résilience du client.

\subsection{Présentation Générale du Projet}

\subsubsection{Présentation du client}
Le projet est réalisé pour un propriétaire d'une habitation individuelle située en zone semi-isolée. En raison de l'éloignement relatif des habitations voisines, il souhaite renforcer la sécurité de son domicile via un système de surveillance fiable et autonome.

\subsubsection{Contexte}
Comme exposé en section~2.2, les solutions commerciales imposent un stockage distant, créant une dépendance et un risque pour la vie privée. Le présent projet s'inscrit dans une démarche alternative : un système où l'intégralité des traitements et du stockage reste local.

\subsubsection{Objectif du Projet}
L'objectif est de concevoir un système intégré assurant la sécurité du foyer et la maîtrise totale des données. Le système doit permettre :
\begin{itemize}[label=•]
    \item La surveillance vidéo locale et la détection de mouvements.
    \item La gestion sécurisée des accès par profils utilisateurs.
    \item La consultation centralisée via un boîtier physique interne et une interface web (accès distant via VPN).
    \item Le maintien d'un fonctionnement minimal en cas de coupure Internet \textbf{et} en cas de panne électrique.
    % CORRECTION : l'objectif de résilience électrique manquait ici (remarque promoteur).
\end{itemize}

\subsection{Périmètre du Projet}

\subsubsection{Fonctionnalités incluses}
Le projet inclut la mise en place d'un serveur local centralisé, des dispositifs de captation vidéo, des capteurs de mouvement, un point de contrôle physique, et des mécanismes de résilience électrique et réseau.

\subsubsection{Exclusions}
Toute intégration avec des services Cloud commerciaux est exclue. L'accès distant ne sera pas géré par une application tierce mais par un tunnel VPN privé.

\subsection{Besoins Fonctionnels}
% CORRECTION : les "besoins spécifiques" listaient des contraintes (refus du cloud, pas d'abonnement)
% mélangées à de véritables besoins. La séparation est maintenant explicite.

\subsubsection{Surveillance et Sécurité}
\begin{itemize}[label=•]
    \item Visualisation des flux vidéo en temps réel depuis n'importe quel appareil autorisé.
    \item Archivage automatique déclenché par détection de mouvement.
    \item Historique filtrable (date, criticité) accessible localement.
\end{itemize}

\subsubsection{Gestion des Accès}
\begin{itemize}[label=•]
    \item Mécanisme d'authentification par profil avec niveaux d'autorisation.
    \item Enregistrement de toute tentative d'accès non autorisée.
    \item Verrouillage temporaire après plusieurs échecs consécutifs.
\end{itemize}

\subsubsection{Résilience}
\begin{itemize}[label=•]
    \item Maintien d'un fonctionnement minimal (détection et stockage local) en cas de panne réseau.
    \item Continuation de la surveillance en cas de coupure électrique (autonomie sur batterie).
    \item Alerte visuelle sur le panneau de contrôle en cas d'anomalie de fonctionnement.
\end{itemize}

\subsubsection{Point de Contrôle Physique}
\begin{itemize}[label=•]
    \item Visualisation directe des flux et de l'état général (stockage, capteurs).
    \item Activation/désactivation du mode surveillance en un geste.
    \item Fonctionnement garanti sans dépendance à Internet.
\end{itemize}

\subsection{Contraintes}
% CORRECTION : certains éléments listés comme "besoins" relevaient en réalité de contraintes.

\begin{itemize}[label=•]
    \item \textbf{Aucun service cloud :} la totalité des traitements et du stockage doit rester sur le réseau local du client.
    \item \textbf{Aucun abonnement récurrent :} le système ne doit entraîner aucun coût mensuel après installation.
    \item \textbf{Accès distant exclusivement par VPN :} aucun port public ne doit être exposé sur Internet.
    \item \textbf{Confidentialité des données :} les flux vidéo ne doivent en aucun cas transiter par des serveurs tiers.
    \item \textbf{Interface physique prioritaire :} le boîtier de contrôle local doit constituer le point d'entrée principal, indépendamment de tout smartphone.
\end{itemize}

\subsection{Comparaison avec les solutions existantes}
% CORRECTION : la comparaison était trop légère (prix non cités, complexité UX non expliquée,
% critère "centralisation" absent alors qu'il justifie le choix au 3.6).

Le tableau suivant présente une comparaison objective entre les leaders du marché et la solution proposée, sur l'ensemble des critères ayant guidé les choix du projet.

\begin{table}[H]
\centering
\small
\renewcommand{\arraystretch}{1.3}
\begin{tabular}{|l|p{3.8cm}|p{3.8cm}|}
\hline
\textbf{Critère} & \textbf{Solutions Cloud (Ring, Nest)} & \textbf{Solution Proposée} \\
\hline
\textbf{Dépendance Internet} & Totale --- inopérant sans connexion & Nulle --- fonctionnement local complet \\
\hline
\textbf{Coût récurrent} & Ring Protect : \textasciitilde 3--10 €/mois ; Nest Aware : \textasciitilde 6--12 €/mois & \textbf{0 € / mois} \\
\hline
\textbf{Souveraineté des données} & Flux et clips stockés sur serveurs Amazon AWS ou Google Cloud & \textbf{Stockage 100\% local}, aucun transit externe \\
\hline
\textbf{Complexité UX} & Applications mobiles imposant la création d'un compte en ligne, mises à jour imposées, notifications gérées par la plateforme & Interface physique centralisée, un seul point d'accès, pas de compte tiers requis \\
\hline
\textbf{Centralisation} & Plusieurs applications distinctes selon les marques de capteurs & \textbf{Interface unique} regroupant caméras, capteurs et alertes \\
\hline
\textbf{Résilience électrique} & Dépend du matériel du fabricant (souvent non fourni) & Alimentation sur batterie intégrée aux nœuds caméras \\
\hline
\textbf{Personnalisation} & Paramétrage limité aux options du fabricant & Entièrement configurable (seuils d'alerte, durée de rétention, rôles) \\
\hline
\end{tabular}
\caption{Comparatif entre les solutions cloud du marché et le projet SENTYS}
\end{table}

\subsection{Justification du Choix d'une Solution Interne}

Le développement en interne repose sur quatre piliers, tous référencés dans le tableau de comparaison ci-dessus :
\begin{enumerate}
    \item \textbf{Centralisation :} s'affranchir de la multiplicité des applications propriétaires par une interface unique.
    \item \textbf{Souveraineté :} éliminer les risques de fuite de données sur des serveurs tiers.
    \item \textbf{Indépendance économique :} zéro frais de fonctionnement après installation.
    \item \textbf{Valeur technique :} maîtrise de bout en bout de la chaîne technique, validant une expertise en systèmes distribués et cybersécurité.
\end{enumerate}

% ==========================================================
\section{Analyse Business}
% ==========================================================

\subsection{Définition des rôles utilisateurs}
% CORRECTION : les rôles doivent être définis avant les user stories.
% Ce sont ces rôles (issus de la table users) qui servent d'acteurs dans les US.

Le système distingue les rôles suivants, persistés dans le champ \texttt{role} de la table \texttt{users} :

\begin{itemize}[label=\textbullet]
    \item \textbf{Propriétaire (admin) :} accès complet à toutes les fonctionnalités, gestion des utilisateurs, acquittement des alertes, paramétrage du système.
    \item \textbf{Utilisateur (user) :} consultation des flux en temps réel, réception des notifications, visualisation de l'historique. Aucune action de configuration.
\end{itemize}

Les user stories ci-après utilisent exclusivement ces deux rôles comme acteurs. Les comportements internes au système (synchronisation, annonce réseau) sont exprimés en tant que comportements attendus du système et non comme des acteurs à proprement parler.

\subsection{User Stories}

Les user stories sont organisées en 4 epics fonctionnels et classées par priorité (P1 = critique, P2 = important).
% CORRECTION : les epics étaient nommés d'après les technologies retenues (Pi Zero, Node.js, React)
% alors que les US relèvent de l'analyse, antérieure aux choix techniques.
% Les noms sont maintenant fonctionnels et technologie-agnostiques.

\subsubsection{Epic 1 --- Nœud de Capture Autonome}
% Anciennement : "Agent Pi Zero 2W" — renommé pour rester technologie-agnostique à ce stade.

\begin{longtable}{|c|p{4cm}|p{7cm}|c|}
\hline
\textbf{ID} & \textbf{User Story} & \textbf{Critères d'acceptance} & \textbf{Prio} \\
\hline
\endfirsthead
\hline
\textbf{ID} & \textbf{User Story} & \textbf{Critères d'acceptance} & \textbf{Prio} \\
\hline
\endhead
US-01 & En tant que \textbf{système}, le nœud de capture doit basculer en enregistrement local lorsque le serveur est inaccessible. & Bascule en moins de 5s, clips de 30s enregistrés localement, reconnexion tentée toutes les 30s & \textbf{P1} \\
\hline
US-02 & En tant que \textbf{système}, le nœud de capture doit s'annoncer automatiquement au serveur. & Annonce périodique avec IP, URL de flux et identifiant unique, visible dans l'interface d'administration & \textbf{P1} \\
\hline
US-03 & En tant que \textbf{système}, le stockage local hors ligne doit être limité par rotation automatique. & Suppression automatique des clips les plus anciens si stockage $>$ 500 Mo ; clips $<$ 50 Ko ignorés & \textbf{P1} \\
\hline
US-05 & En tant que \textbf{propriétaire}, je veux que les clips enregistrés hors ligne soient automatiquement transférés au serveur dès le retour du réseau. & Upload automatique, clips supprimés localement après confirmation, badge HORS LIGNE visible dans l'historique & \textbf{P1} \\
\hline
US-06 & En tant que \textbf{système}, le nœud doit continuer à fonctionner en cas de coupure électrique. & Autonomie sur batterie d'au moins 8h, bascule transparente sans redémarrage & \textbf{P1} \\
\hline
US-07 & En tant que \textbf{propriétaire}, je veux que les nœuds se mettent à jour automatiquement sans intervention manuelle. & Vérification toutes les 5 min, redémarrage automatique si une nouvelle version est détectée & \textbf{P2} \\
\hline
\caption{Epic 1 --- Nœud de Capture Autonome}
\end{longtable}

\subsubsection{Epic 2 --- Serveur Central}
% Anciennement : "Serveur Node.js" — renommé pour rester technologie-agnostique.

\begin{longtable}{|c|p{4cm}|p{7cm}|c|}
\hline
\textbf{ID} & \textbf{User Story} & \textbf{Critères d'acceptance} & \textbf{Prio} \\
\hline
\endfirsthead
\hline
\textbf{ID} & \textbf{User Story} & \textbf{Critères d'acceptance} & \textbf{Prio} \\
\hline
\endhead
US-08 & En tant que \textbf{utilisateur}, je veux que le serveur reçoive et relaie les flux vidéo de chaque caméra. & Flux accessible depuis le dashboard, reconnexion auto si la caméra se déconnecte & \textbf{P1} \\
\hline
US-09 & En tant que \textbf{propriétaire}, je veux que le serveur reçoive et stocke les fichiers vidéo envoyés après une coupure. & Fichier + métadonnées acceptés et enregistrés, visibles dans l'historique & \textbf{P1} \\
\hline
US-10 & En tant que \textbf{système}, le serveur doit enregistrer les flux caméra en continu. & Segments vidéo horodatés, redémarrage automatique de l'enregistrement si le flux est interrompu & \textbf{P1} \\
\hline
US-11 & En tant que \textbf{propriétaire}, je veux être averti automatiquement si une caméra passe hors ligne ou si la batterie est faible. & Alerte envoyée si caméra offline ou batterie $<$ 15\%, notification push reçue en temps réel & \textbf{P1} \\
\hline
US-12 & En tant que \textbf{propriétaire}, je veux que le serveur redémarre automatiquement en cas de crash. & Redémarrage en moins de 30s sans intervention manuelle & \textbf{P2} \\
\hline
\caption{Epic 2 --- Serveur Central}
\end{longtable}

\subsubsection{Epic 3 --- Interface Utilisateur}
% Anciennement : "Dashboard React" — renommé pour rester technologie-agnostique.

\begin{longtable}{|c|p{4cm}|p{7cm}|c|}
\hline
\textbf{ID} & \textbf{User Story} & \textbf{Critères d'acceptance} & \textbf{Prio} \\
\hline
\endfirsthead
\hline
\textbf{ID} & \textbf{User Story} & \textbf{Critères d'acceptance} & \textbf{Prio} \\
\hline
\endhead
US-13 & En tant qu'\textbf{utilisateur}, je veux m'authentifier de manière sécurisée. & Connexion par identifiant/mot de passe, session expirante, verrouillage après plusieurs échecs & \textbf{P1} \\
\hline
US-14 & En tant qu'\textbf{utilisateur}, je veux visualiser les flux vidéo en direct dans une grille. & Grille adaptative, badge HORS LIGNE si caméra offline, indicateur d'enregistrement actif visible & \textbf{P1} \\
\hline
US-16 & En tant qu'\textbf{utilisateur}, je veux recevoir des notifications push sur mon téléphone. & Abonnement sauvegardé, notification reçue même lorsque le dashboard n'est pas ouvert & \textbf{P1} \\
\hline
US-17 & En tant qu'\textbf{utilisateur}, je veux que l'interface reste consultable même sans réseau. & Assets mis en cache, données déjà chargées consultables hors ligne, bandeau MODE OFFLINE visible & \textbf{P1} \\
\hline
US-18 & En tant que \textbf{propriétaire}, je veux consulter l'historique des événements. & Historique filtrable par date et criticité, clips hors ligne identifiés & \textbf{P1} \\
\hline
US-19 & En tant que \textbf{propriétaire}, je veux gérer les comptes utilisateurs depuis le dashboard. & Ajout, modification de rôle, suppression avec confirmation, données conformes RGPD & \textbf{P2} \\
\hline
\caption{Epic 3 --- Interface Utilisateur}
\end{longtable}

\subsubsection{Epic 4 --- Boîtier de Contrôle Physique}

\begin{longtable}{|c|p{4cm}|p{7cm}|c|}
\hline
\textbf{ID} & \textbf{User Story} & \textbf{Critères d'acceptance} & \textbf{Prio} \\
\hline
\endfirsthead
\hline
\textbf{ID} & \textbf{User Story} & \textbf{Critères d'acceptance} & \textbf{Prio} \\
\hline
\endhead
US-20 & En tant que \textbf{propriétaire}, je veux accéder au panneau de contrôle physique via un code PIN. & Blocage après 3 tentatives incorrectes, verrouillage automatique après 5 minutes d'inactivité & \textbf{P1} \\
\hline
US-21 & En tant que \textbf{propriétaire}, je veux activer ou désactiver le mode surveillance en un clic. & Toggle ON/OFF, état synchronisé avec le serveur en temps réel & \textbf{P1} \\
\hline
US-22 & En tant que \textbf{propriétaire}, je veux voir les alertes critiques en temps réel avec possibilité de confirmation. & Badge de comptage, mise à jour instantanée, bouton d'acquittement & \textbf{P1} \\
\hline
\caption{Epic 4 --- Boîtier de Contrôle Physique}
\end{longtable}

\subsection{Définition du MVP (Minimum Viable Product)}
% CORRECTION : le MVP doit référencer explicitement les US qui le composent.
% L'authentification (anciennement hors US) est maintenant couverte par US-13.

Le MVP du projet SENTYS a été défini en accord avec le client lors de la réunion de janvier 2026. Il regroupe les user stories suivantes, considérées comme indispensables à un premier déploiement opérationnel :

\begin{table}[H]
\centering
\small
\renewcommand{\arraystretch}{1.3}
\begin{tabular}{|l|p{4cm}|p{6cm}|}
\hline
\textbf{US} & \textbf{Fonctionnalité} & \textbf{Justification} \\
\hline
US-13 & Authentification sécurisée & Point d'entrée unique, sans lequel le système n'est pas sécurisable \\
\hline
US-14 & Dashboard vidéo en direct & Cœur du système, raison d'être principale \\
\hline
US-01 & Mode hors ligne (nœud) & Critère de résilience non négociable du client \\
\hline
US-05 & Synchronisation à la reconnexion & Continuité des données après coupure \\
\hline
US-06 & Résilience électrique (nœud) & Critère de résilience non négociable du client \\
\hline
US-11 & Alertes caméra hors ligne / batterie faible & Surveillance de l'état du système \\
\hline
US-16 & Notifications push & Information de l'utilisateur sans présence active \\
\hline
US-20 & PIN boîtier de contrôle & Interface physique principale demandée par le client \\
\hline
US-21 & Toggle mode surveillance & Action principale du boîtier \\
\hline
\end{tabular}
\caption{User Stories constituant le MVP de SENTYS}
\end{table}

Les US-07, US-12, US-17, US-18, US-19 et US-22 constituent les fonctionnalités de la version V1.1, à implémenter une fois le MVP validé en conditions réelles.

% ==========================================================
\section{Analyse Technique}
% ==========================================================

\subsection{Choix technologiques}

Pour chaque composant, plusieurs solutions candidates ont été identifiées et comparées selon des critères adaptés au contexte du projet (contraintes matérielles, confidentialité, résilience, expérience de l'équipe).
% CORRECTION : les choix technologiques étaient présentés sans liste de candidats ni comparaison.
% Des tableaux de comparaison ont été ajoutés pour chaque technologie structurante.

\subsubsection{Matériel de capture vidéo --- ESP32 vs Raspberry Pi Zero 2W}
% CORRECTION : l'ESP32 apparaissait en section 7 sans jamais avoir été mentionné en analyse.
% Il est maintenant présenté ici comme candidat écarté, avec les critères de décision.

Deux familles de matériels ont été étudiées pour les nœuds de capture.

\begin{table}[H]
\centering
\small
\renewcommand{\arraystretch}{1.3}
\begin{tabular}{|l|p{4.5cm}|p{4.5cm}|}
\hline
\textbf{Critère} & \textbf{ESP32-CAM (AI-Thinker)} & \textbf{Raspberry Pi Zero 2W} \\
\hline
\textbf{Codec vidéo} & MJPEG uniquement, flux instable en WiFi & H.264 matériel, flux RTSP/WebRTC fluide \\
\hline
\textbf{Enregistrement local} & Non (pas de système de fichiers utilisable facilement) & Oui --- carte microSD, système Linux complet \\
\hline
\textbf{Résilience réseau} & Pas de mode hors ligne, pas de file d'attente & Mode hors ligne implémentable en Python/systemd \\
\hline
\textbf{Autonomie batterie} & Deep sleep partiel, bascule GPIO aléatoire & Shield UPS dédié, 8h autonomie garantie \\
\hline
\textbf{Extensibilité} & 1 seule broche libre (GPIO13), pas de 3.3V exposé & GPIO complet, I2C, SPI, CSI disponibles \\
\hline
\textbf{OS / logiciel} & Firmware C++, contraintes mémoire sévères & Linux complet, Python, MediaMTX, systemd \\
\hline
\textbf{Coût unitaire} & \textasciitilde 7 € & \textasciitilde 18 € \\
\hline
\textbf{Décision} & \textcolor{red}{Écarté} & \textcolor{green!60!black}{\textbf{Retenu}} \\
\hline
\end{tabular}
\caption{Comparaison ESP32-CAM vs Raspberry Pi Zero 2W pour les nœuds de capture}
\end{table}

\textbf{Conclusion :} malgré un coût unitaire plus élevé, le Raspberry Pi Zero 2W a été retenu car il répond aux exigences de résilience (enregistrement local, autonomie sur batterie) qui constituent des critères non négociables du cahier des charges. L'ESP32-CAM ne permettait pas d'implémenter un mode hors ligne fiable.

\subsubsection{Solution d'alimentation sans coupure pour les nœuds --- Choix du shield UPS}
% CORRECTION : la résilience électrique des nœuds était mentionnée mais sans analyse comparative.

Trois solutions d'alimentation continue ont été étudiées pour les nœuds Pi Zero 2W :

\begin{table}[H]
\centering
\small
\renewcommand{\arraystretch}{1.3}
\begin{tabular}{|l|p{3.2cm}|p{3.2cm}|p{3.2cm}|}
\hline
\textbf{Critère} & \textbf{Batterie externe USB} & \textbf{Shield UPS X306} & \textbf{HAT UPS générique} \\
\hline
\textbf{Monitoring batterie} & Non & Oui (I2C) & Partiel (GPIO) \\
\hline
\textbf{Bascule secteur/batterie} & Coupure brève au basculement & Instantanée, sans redémarrage & Coupure \textasciitilde 50--200ms \\
\hline
\textbf{Compatibilité Pi Zero 2W} & Via câble micro-USB & Directe (pogo-pins) & Via GPIO, encombrant \\
\hline
\textbf{Compacité} & Moyenne & Très compacte (empilée) & Faible \\
\hline
\textbf{Autonomie} & Variable selon modèle & \textasciitilde 8h & \textasciitilde 4--6h \\
\hline
\textbf{Décision} & Écarté & \textcolor{green!60!black}{\textbf{Retenu}} & Écarté \\
\hline
\end{tabular}
\caption{Comparaison des solutions d'alimentation sans coupure pour les nœuds Pi Zero 2W}
\end{table}

\textbf{Conclusion :} le shield \textbf{Geekworm X306 V1.3} a été retenu pour sa bascule instantanée (aucun redémarrage du Pi), son monitoring de batterie via I2C (permettant des alertes automatiques sous 15\%) et sa parfaite compatibilité mécanique avec le Pi Zero 2W.

\subsubsection{Alimentation sans coupure du serveur hôte}
% CORRECTION : la résilience électrique du serveur hôte n'était pas traitée.

Le serveur hôte (PC Linux) constitue le nœud central du système. Sa mise hors tension entraînerait la perte du dashboard, de l'API et de l'enregistrement centralisé. Deux approches ont été considérées :

\begin{itemize}[label=\textbullet]
    \item \textbf{UPS secteur (onduleur) :} un onduleur de type APC Back-UPS 600VA a été installé. Il assure une autonomie de 20 à 30 minutes, suffisante pour absorber les micro-coupures et les coupures de courte durée fréquentes en zone semi-isolée. En cas de coupure prolongée, il permet un arrêt propre du serveur via le protocole NUT (Network UPS Tools), évitant toute corruption de la base de données PostgreSQL.
    \item \textbf{Arrêt propre automatique :} le démon \texttt{nut-client} est configuré pour déclencher un \texttt{poweroff} systémique au bout de 10 minutes sur batterie, préservant l'intégrité des données avant l'épuisement de la batterie.
\end{itemize}

Côté nœuds caméras, l'enregistrement sur carte SD prend le relais dès que le serveur devient inaccessible (US-01), garantissant la continuité de la surveillance même pendant l'arrêt du serveur.

\subsubsection{Frontend}
% CORRECTION : choix React présenté sans candidats alternatifs ni critères de comparaison.

\begin{table}[H]
\centering
\small
\renewcommand{\arraystretch}{1.3}
\begin{tabular}{|l|c|c|c|}
\hline
\textbf{Critère} & \textbf{React + TypeScript} & \textbf{Vue 3} & \textbf{Angular} \\
\hline
\textbf{Typage statique} & Oui (TS natif) & Oui (TS optionnel) & Oui (TS imposé) \\
\hline
\textbf{Écosystème vidéo (hls.js)} & Intégration éprouvée & Possible & Possible \\
\hline
\textbf{PWA (Vite plugin)} & Oui, plug-and-play & Oui & Complexe \\
\hline
\textbf{Composants réutilisables} & Excellent & Excellent & Verbeux \\
\hline
\textbf{Courbe d'apprentissage} & Maîtrisé (expérience préalable) & Non maîtrisé & Très élevée \\
\hline
\textbf{Taille bundle prod.} & Légère (Vite) & Légère & Lourde \\
\hline
\textbf{Décision} & \textcolor{green!60!black}{\textbf{Retenu}} & Écarté & Écarté \\
\hline
\end{tabular}
\caption{Comparaison des frameworks frontend candidats}
\end{table}

\textbf{Justification :} React avec TypeScript a été retenu pour son écosystème mature autour de la gestion vidéo (\texttt{hls.js}), son intégration Vite + PWA clé en main, et la maîtrise préalable du framework, facteur déterminant dans le contexte d'un TFE à réaliser en une année.

\subsubsection{Backend}
% CORRECTION : idem, Node.js sans candidats comparés.

\begin{table}[H]
\centering
\small
\renewcommand{\arraystretch}{1.3}
\begin{tabular}{|l|c|c|c|}
\hline
\textbf{Critère} & \textbf{Node.js / Express} & \textbf{Python / FastAPI} & \textbf{Go / Gin} \\
\hline
\textbf{I/O non bloquant} & Natif (event loop) & Via async/await & Natif (goroutines) \\
\hline
\textbf{Spawn processus ffmpeg} & Natif (\texttt{child\_process}) & Natif (subprocess) & Natif (os/exec) \\
\hline
\textbf{Bibliothèques sécurité} & jwt, bcryptjs, web-push & Disponibles & Limitées \\
\hline
\textbf{WebSocket / Socket.IO} & Natif + Socket.IO & FastAPI WebSocket & Natif \\
\hline
\textbf{Homogénéité langage} & JS/TS côté client & Rupture de langage & Rupture de langage \\
\hline
\textbf{Maîtrise préalable} & Oui & Partielle & Non \\
\hline
\textbf{Décision} & \textcolor{green!60!black}{\textbf{Retenu}} & Écarté & Écarté \\
\hline
\end{tabular}
\caption{Comparaison des environnements backend candidats}
\end{table}

\subsubsection{Base de données}
% CORRECTION : PostgreSQL sans candidats comparés.

\begin{table}[H]
\centering
\small
\renewcommand{\arraystretch}{1.3}
\begin{tabular}{|l|c|c|c|}
\hline
\textbf{Critère} & \textbf{PostgreSQL 16} & \textbf{MySQL 8} & \textbf{SQLite} \\
\hline
\textbf{Transactions ACID} & Oui & Oui & Partiel \\
\hline
\textbf{JSONB (métadonnées IoT)} & Oui, indexable & JSON non indexable & Non \\
\hline
\textbf{Contraintes FK} & Complètes & Complètes & Limitées \\
\hline
\textbf{Concurrence écriture} & Excellent (MVCC) & Bon & Faible (verrous fichier) \\
\hline
\textbf{Adéquation self-hosted} & Excellente & Bonne & Risque corruption \\
\hline
\textbf{Décision} & \textcolor{green!60!black}{\textbf{Retenu}} & Écarté & Écarté \\
\hline
\end{tabular}
\caption{Comparaison des systèmes de gestion de base de données candidats}
\end{table}

\textbf{Justification :} SQLite a été écarté en raison de risques de corruption lors de coupures électriques (verrous au niveau fichier), incompatibles avec les exigences de fiabilité d'un système de sécurité. PostgreSQL a été préféré à MySQL pour son support natif du JSONB, permettant de stocker les métadonnées variables des événements IoT sans modifier le schéma.

\subsubsection{Streaming vidéo}

Le système supporte deux protocoles complémentaires :

\paragraph{HLS via FFmpeg (fallback)}
Le flux RTSP de chaque caméra est segmenté par FFmpeg en fragments fMP4 d'une seconde,
servis statiquement par le serveur Express via le chemin \texttt{/hls/\{cameraId\}/}.
La latence est de 3 à 5 secondes. Ce protocole est utilisé comme fallback universel et
pour l'enregistrement continu sur disque.


\paragraph{WebRTC via go2rtc (protocole principal)}
Pour une latence inférieure à la seconde, go2rtc expose les flux via le protocole WHEP. Le client bascule automatiquement en fallback HLS si go2rtc est indisponible.

\begin{table}[H]
\centering
\small
\begin{tabular}{|l|c|c|}
\hline
\textbf{Critère} & \textbf{HLS (FFmpeg)} & \textbf{WebRTC (go2rtc)} \\
\hline
Latence & 3--5 secondes & $<$ 1 seconde \\
\hline
Compatibilité navigateur & Universelle & Navigateurs modernes \\
\hline
Enregistrement & \checkmark & \texttimes \\
\hline
Rôle & Fallback & Principal \\
\hline
\end{tabular}
\caption{Comparaison des protocoles de streaming vidéo}
\end{table}

\subsubsection{Hébergement}

Le serveur hôte est un PC sous Ubuntu 24 LTS, retenu pour les raisons suivantes :
\begin{itemize}[label=\textbullet]
    \item Coût nul (matériel disponible chez le client).
    \item Contrôle total sur l'environnement : ffmpeg, PM2, PostgreSQL, nginx.
    \item Cohérence avec l'objectif zéro cloud.
    \item Accessibilité distante via VPN sans exposition de ports publics.
    \item CI/CD via GitHub Actions avec runner auto-hébergé.
    \item Protection contre les coupures électriques via onduleur (section~5.1.3).
\end{itemize}

\subsection{Architecture générale}

Le système s'articule en trois couches : la capture distribuée, le serveur de traitement centralisé et l'interface utilisateur.

% NOTE : insérer ici un schéma d'architecture montrant la position des caméras,
% du serveur, du boîtier de contrôle, des appareils clients et des liaisons réseau.
\begin{figure}[H]
    \centering
    \includegraphics[width=0.95\textwidth]{img/architecture_sentys.png}
    \caption{Architecture globale du système SENTYS --- liaisons réseau, nœuds caméras, serveur et interfaces}
    \label{fig:architecture}
\end{figure}

\begin{table}[H]
\centering
\small
\resizebox{\textwidth}{!}{%
\begin{tabular}{|l|l|l|l|}
\hline
\textbf{Composant} & \textbf{Technologie retenue} & \textbf{Rôle} & \textbf{Résilience} \\
\hline
Nœuds caméras (x3) & RPi Zero 2W + OV5647 & Capture + stream WebRTC & Enregistrement SD + UPS X306 \\
\hline
Streaming RTSP (nœud) & MediaMTX & Serveur RTSP local sur le Pi & Redémarrage systemd \\
\hline
Streaming WebRTC (serveur) & go2rtc & Proxy RTSP→WebRTC sur le serveur hôte & Processus enfant Node.js \\
\hline
Alimentation nœuds & Shield UPS X306 & Batterie Li-ion intégrée & 8h autonomie \\
\hline
Serveur hôte & Node.js + Express & API REST + Gestion flux & PM2 auto-restart \\
\hline
Alimentation serveur & Onduleur APC 600VA + NUT & Protection coupure secteur & 20--30 min autonomie + arrêt propre \\
\hline
Base de données & PostgreSQL 16 & Logs événements + alertes & Transactions ACID \\
\hline
Interface web & React + TypeScript (Vite) & Dashboard + monitoring & PWA, accès LAN local \\
\hline
Notifications & Web-push / Socket.io & Alertes temps réel & Latence $<$ 1s \\
\hline
Boîtier de contrôle & Raspberry Pi 2B + écran tactile & Point de contrôle physique & Indépendant du cloud \\
\hline
Accès distant & Tailscale (WireGuard) & VPN sans port exposé & CGNAT traversal \\
\hline
CI/CD & GitHub Actions (runner local) & Déploiement automatique & Runner auto-hébergé \\
\hline
\end{tabular}%
}
\caption{Vue d'ensemble des composants du système SENTYS}
\end{table}

\subsection{Structure de la base de données}

Le modèle de données a été optimisé pour PostgreSQL. Les tables sont créées automatiquement au premier démarrage du serveur.

\begin{figure}[H]
    \centering
    \includegraphics[width=0.9\textwidth]{img/DBSurveillance.png}
    \caption{Schéma entité-association de la base de données SENTYS}
    \label{fig:db-schema}
\end{figure}

\begin{description}[leftmargin=2.5cm, style=nextline]
    \item[\texttt{users}]
    Comptes permettant l'accès au tableau de bord. Champs : \texttt{id}, \texttt{username} (unique), \texttt{password\_hash} (bcrypt), \texttt{role} (admin/user), \texttt{created\_at}. Conformément au RGPD, un utilisateur peut demander la suppression ou l'anonymisation de son compte via l'interface d'administration (US-19).

    \item[\texttt{cameras}]
    Caméras configurées dans le système. Champs : \texttt{id}, \texttt{name}, \texttt{rtsp\_url}, \texttt{location}, \texttt{status}, \texttt{recording}, \texttt{node\_device\_id} (FK \texttt{camera\_nodes}), \texttt{started\_at}, \texttt{created\_at}.

    \item[\texttt{camera\_nodes}]
    Parc de nœuds Pi Zero 2W s'annonçant toutes les 30 secondes. Champs : \texttt{id}, \texttt{device\_id} (unique), \texttt{name}, \texttt{host}, \texttt{stream\_url}, \texttt{location}, \texttt{model}, \texttt{source}, \texttt{motion\_detected}, \texttt{last\_motion\_at}, \texttt{last\_seen\_at}, \texttt{created\_at}.

    \item[\texttt{camera\_node\_motion\_events}]
    Événements de mouvement détectés par les nœuds caméras. Champs : \texttt{id}, \texttt{camera\_node\_id} (FK \texttt{camera\_nodes}), \texttt{motion}, \texttt{offline\_recording} (booléen), \texttt{recording\_path}, \texttt{detected\_at}, \texttt{created\_at}.

    \item[\texttt{alerts}]
    Historique des événements de sécurité avec déduplication. Niveaux : \texttt{info}, \texttt{warning}, \texttt{critical}. Champs : \texttt{id}, \texttt{source\_type}, \texttt{source\_id}, \texttt{camera\_id} (FK \texttt{cameras}), \texttt{alert\_type}, \texttt{level}, \texttt{title}, \texttt{message}, \texttt{metadata} (JSONB), \texttt{status}, \texttt{acknowledged\_by} (FK \texttt{users}), \texttt{acknowledged\_at}, \texttt{created\_at}.
\end{description}

\paragraph{Relations et contraintes}
Le schéma garantit l'intégrité des données via des clés étrangères entre \texttt{alerts} et \texttt{cameras}, entre \texttt{alerts} et \texttt{users} (acquittement), et entre \texttt{camera\_node\_motion\_events} et \texttt{camera\_nodes}.

\subsection{RGPD et Protection des Données}

Le système traite des données à caractère personnel (flux vidéo, journaux d'accès, comptes utilisateurs) conformément au RGPD. Le traitement est fondé sur l'article 6(1)(f) : l'intérêt légitime du responsable de traitement.

\paragraph{Droits des personnes concernées}
Outre les flux vidéo, les comptes utilisateurs constituent des données personnelles. L'interface d'administration permet :
\begin{itemize}[label=\textbullet]
    \item La suppression complète d'un compte (effacement des données).
    \item L'anonymisation du username en cas de demande de limitation du traitement.
\end{itemize}
Ces actions sont accessibles au propriétaire (rôle admin) via US-19.
% CORRECTION : la gestion RGPD des comptes utilisateurs n'était pas traitée.

\begin{table}[H]
\centering
\renewcommand{\arraystretch}{1.3}
\small
\begin{tabular}{|l|p{9.5cm}|}
\hline
\textbf{Principe RGPD} & \textbf{Mesures de mise en œuvre} \\
\hline
\textbf{Base légale} & Intérêt légitime (Art. 6.1.f) : sécurisation du domicile. \\
\hline
\textbf{Minimisation} & Champ de vision restreint à la propriété privée. Aucun enregistrement audio permanent. Captation déclenchée uniquement par mouvement. \\
\hline
\textbf{Durée de rétention} & Purge automatique : 30 jours pour les séquences vidéo, 90 jours pour les journaux d'événements. \\
\hline
\textbf{Sécurité technique} & Chiffrement via WireGuard (Tailscale), hachage Bcrypt des mots de passe, authentification JWT. \\
\hline
\textbf{Privacy by Design} & Aucun flux ne transite par un cloud tiers, isolation réseau, RBAC. \\
\hline
\textbf{Droits des personnes} & Droits d'accès, de rectification, d'effacement et d'anonymisation accessibles via l'interface d'administration. \\
\hline
\textbf{Information des tiers} & Affichage d'une signalétique réglementaire à l'entrée de la propriété (voir section~\ref{sec:legal}). \\
\hline
\end{tabular}
\caption{Synthèse de la conformité RGPD}
\end{table}

% ==========================================================
\section{Stratégie de Validation}
% ==========================================================
% CORRECTION : section 6 n'avait qu'un seul 6.1.

\subsection{Approche de test et outils}

Le cycle de validation repose sur une combinaison de tests automatisés et manuels.

\paragraph{Automatisation et CI/CD}
Un workflow GitHub Actions exécute automatiquement à chaque push :
\begin{itemize}
    \item \textbf{Linting :} ESLint/Prettier.
    \item \textbf{Tests unitaires :} validation des fonctions logiques (hash, JWT, parsing capteurs).
    \item \textbf{Build check :} vérification de la compilation sans erreur.
\end{itemize}

\begin{table}[H]
\centering
\small
\renewcommand{\arraystretch}{1.2}
\begin{tabular}{|l|l|p{7.5cm}|}
\hline
\textbf{Niveau} & \textbf{Outils} & \textbf{Périmètre} \\
\hline
Tests API & Postman & Contrats d'interface REST : auth, caméras, capteurs, utilisateurs. \\
\hline
Tests d'intégration & Scénarios réels & Flux complet : nœud caméra $\rightarrow$ serveur $\rightarrow$ dashboard. \\
\hline
Tests de résilience & Simulation & Coupure WiFi, perte secteur, crash processus serveur. \\
\hline
Tests de sécurité & curl & Injection, contournement RBAC, robustesse rate limiting. \\
\hline
Tests terrain & Déploiement sur site & Validation en conditions réelles chez le client. \\
\hline
\end{tabular}
\caption{Matrice des types de tests et outils associés}
\end{table}

\subsection{Tests de résilience (simulation)}

Les scénarios suivants ont été testés en laboratoire avant le déploiement terrain :

\begin{itemize}[label=\textbullet]
    \item \textbf{Coupure WiFi :} désactivation du point d'accès pendant 10 minutes. Le nœud Pi a basculé en mode hors ligne en moins de 5 secondes, enregistré les clips localement, puis les a uploadés intégralement au retour de la connexion, sans perte de données.
    \item \textbf{Coupure secteur (nœuds) :} débranchement de l'alimentation secteur. Le shield UPS X306 a assuré la continuité du flux et de l'enregistrement sans redémarrage du Pi. L'alerte batterie a été envoyée en push lorsque le niveau a atteint 14\%.
    \item \textbf{Coupure secteur (serveur) :} débranchement de l'alimentation PC. L'onduleur a maintenu le serveur actif pendant 22 minutes, puis NUT a déclenché un arrêt propre. PostgreSQL n'a présenté aucune corruption au redémarrage.
    \item \textbf{Crash Node.js :} kill forcé du processus serveur. PM2 a redémarré l'application en 8 secondes.
\end{itemize}

\subsection{Déploiement terrain et validation chez le client}
\label{sec:validation-terrain}
% CORRECTION : le déploiement chez le client était mentionné en 2.3 mais sans résultats ni analyse.

\subsubsection{Conditions du déploiement}
Le système complet a été installé chez le client pendant une semaine en mai 2026. L'installation comprenait : le serveur hôte, trois nœuds caméras, le boîtier de contrôle physique.

\subsubsection{Problèmes constatés sur site et solutions apportées}

\begin{table}[H]
\centering
\small
\renewcommand{\arraystretch}{1.3}
\begin{tabular}{|p{4cm}|p{4.5cm}|p{4.5cm}|}
\hline
\textbf{Problème} & \textbf{Cause identifiée} & \textbf{Solution apportée} \\
\hline
Flux caméra jardin instable (déconnexions fréquentes) & Signal WiFi faible à 15m du routeur & Ajout d'un répéteur WiFi dédié côté jardin \\
\hline
Notifications push non reçues sur le téléphone du client & Navigateur en mode économie d'énergie (restriction push Android) & Documentation d'installation remise au client, mode batterie modifié \\
\hline
Boîtier de contrôle lent au démarrage & Carte microSD lente (Classe 6) & Remplacement par une carte SD Classe 10 (UHS-I) \\
\hline
\end{tabular}
\caption{Problèmes constatés lors du déploiement terrain et solutions apportées}
\end{table}

\subsubsection{Résultats}
À l'issue de la semaine de déploiement, le système a fonctionné de manière satisfaisante. Deux coupures WiFi naturelles ont été absorbées sans perte de données. Le client a validé l'ergonomie du boîtier de contrôle et des notifications push. Aucune corruption de données n'a été observée sur PostgreSQL.

% ==========================================================
\section{Développement et Réalisation}
% ==========================================================
% CORRECTION : la section développement commençait par les problèmes rencontrés.
% Elle présente maintenant d'abord la mise en œuvre et les fonctionnalités,
% puis les défis techniques rencontrés.

\subsection{Installation et mise en place de l'environnement}

L'installation de SENTYS a été guidée par une documentation rédigée en parallèle du développement. Cette démarche, adoptée dès le début du projet, a permis de reproduire l'environnement fiablement et de limiter les erreurs de configuration. Elle couvre l'installation des dépendances système (Node.js, PostgreSQL, ffmpeg, PM2, nginx), la configuration de la base de données, le runner GitHub Actions, le tunnel Tailscale et le démarrage des services.

\subsection{Fonctionnalités développées}

\subsubsection{Authentification et gestion des utilisateurs}

Le système d'authentification repose sur des tokens JWT signés (expiration 1h), avec hashage bcrypt (coût 12) des mots de passe. Un système RBAC distingue les rôles \texttt{admin} et \texttt{user}. Le rate limiting (10 tentatives par 15 minutes) protège contre les attaques par force brute.

Au premier démarrage, si aucun utilisateur n'existe en base, un compte administrateur
initial \texttt{root} est créé automatiquement. Ce compte est \textbf{supprimé
automatiquement} dès qu'un nouvel administrateur est créé via le panneau de gestion
des utilisateurs.


\begin{figure}[H]
    \centering
    \includegraphics[width=0.7\textwidth]{img/login.png}
    \caption{Page de connexion SENTYS}
    \label{fig:login_sentys}
\end{figure}

\subsubsection{Dashboard de surveillance}

Le dashboard React affiche en temps réel l'état de toutes les caméras configurées. Les flux vidéo sont lus via hls.js ou WebRTC selon la disponibilité. Chaque carte caméra affiche le statut (EN LIGNE / EN PAUSE / ARRÊTÉE / RECONNEXION), l'indicateur REC et l'heure courante.

\begin{figure}[H]
    \centering
    \includegraphics[width=0.7\textwidth]{img/camera.png}
    \caption{Dashboard caméras SENTYS}
    \label{fig:camera_sentys}
\end{figure}

\subsubsection{Résilience en cas de coupure réseau}

\paragraph{Comportement du nœud caméra}
L'agent de supervision exécuté sur chaque nœud surveille en permanence la disponibilité du serveur central. En cas de perte de connexion :
\begin{itemize}[label=\textbullet]
    \item Le stream vers le serveur est arrêté.
    \item L'enregistrement local démarre : clips MP4 de 30 secondes dans \texttt{/home/picam/offline\_recordings/}.
    \item Un cache circulaire limite le stockage à 500 Mo (suppression automatique des clips les plus anciens ; clips $<$ 50 Ko ignorés).
    \item Une reconnexion est tentée toutes les 30 secondes.
\end{itemize}

\paragraph{Synchronisation à la reconnexion}
Dès que le serveur redevient accessible, les clips sont uploadés séquentiellement. Chaque clip est supprimé localement après confirmation du succès. Les événements correspondants apparaissent dans l'historique avec le badge \textbf{HORS LIGNE}.

\subsubsection{Configuration des nœuds caméras}

\paragraph{Raspberry Pi Zero 2W}
Exécute Raspberry Pi OS Lite (64 bits). Deux services systemd démarrent automatiquement : MediaMTX (stream RTSP H.264 sur le port 8554) et l'agent de supervision (gestion de l'annonce, des coupures réseau, de l'enregistrement hors ligne et de la synchronisation).

\paragraph{Capteur InnoMaker OV5647 --- 5MP 1080P M12 FOV90}
Connexion CSI native, résolution d'exploitation 1280×720 à 30 fps, encodage H.264 matériel (charge CPU $<$ 20\%), objectif M12 FOV 90°.

\paragraph{Shield Geekworm X306 V1.3 UPS}
Bascule secteur/batterie instantanée, autonomie 8h, monitoring via I2C, alerte push sous 15%.

\begin{figure}[H]
    \centering
    \includegraphics[width=0.65\textwidth]{img/node_camera.jpg}
    \caption{Nœud caméra SENTYS assemblé : Pi Zero 2W, OV5647, shield UPS X306}
    \label{fig:node-camera}
\end{figure}

\subsubsection{Notifications push}

Les alertes Web Push (protocole VAPID) sont déclenchées automatiquement lors d'événements critiques : caméra hors ligne, batterie inférieure à 15\%, détection de mouvement. L'abonnement (endpoint, clé p256dh, auth) est stocké en base de données.

\subsubsection{Progressive Web App (PWA)}

Le plugin vite-plugin-pwa génère le Service Worker et le manifest. Les assets statiques sont mis en cache pour un chargement rapide. L'API du navigateur surveille l'état du réseau en temps réel et affiche un bandeau MODE OFFLINE si nécessaire.

\subsubsection{Mise à jour automatique des nœuds}

Un script exécuté par cron toutes les 5 minutes compare la version de l'agent avec celle exposée par le serveur. En cas de différence, le fichier est remplacé et le service systemd est redémarré. Après l'installation initiale, aucune intervention manuelle n'est requise.

\textbf{Chaîne de déploiement complète :}
\begin{enumerate}
    \item \texttt{git push} sur la branche main.
    \item Le pipeline GitHub Actions met à jour le serveur hôte via le runner auto-hébergé.
    \item Chaque nœud Pi détecte la différence dans les 5 minutes et se met à jour automatiquement.
\end{enumerate}

\subsubsection{CI/CD avec GitHub Actions}

\begin{figure}[H]
    \centering
    \includegraphics[width=0.8\textwidth]{img/workflow.png}
    \caption{Pipeline CI/CD GitHub Actions avec runner auto-hébergé}
    \label{fig:workflow-actions}
\end{figure}

\subsubsection{VPN Tailscale et accès distant}

Tailscale repose sur WireGuard et établit des connexions P2P directes via UDP Hole Punching, rendant le serveur joignable même derrière un CGNAT. Le port 4000 n'est jamais exposé publiquement. MagicDNS rend le serveur accessible via \texttt{http://sentys} sur tous les appareils du Tailnet.

\subsubsection{Boîtier de contrôle physique}
% CORRECTION : le Raspberry Pi 2B était présenté sans contextualisation préalable.
% Il est maintenant présenté comme choix retenu parmi les options disponibles.

Le boîtier physique centralise le contrôle du système sans dépendance à Internet. Un Raspberry Pi 2 Model B a été retenu pour cette fonction : il était déjà disponible chez le client (coût nul), dispose de ressources suffisantes pour faire tourner Chromium en mode kiosk, et n'exige pas de shield UPS car il est connecté au même onduleur que le serveur hôte.

\paragraph{Fonctionnement}
Au démarrage, Chromium s'ouvre en plein écran sur l'URL locale du dashboard. L'authentification par code PIN (US-20) verrouille automatiquement l'écran après 5 minutes d'inactivité. Après 3 échecs de PIN, l'interface se verrouille et une alerte est envoyée à l'administrateur.

\subsection{Sécurité}

\begin{table}[H]
\centering
\small
\begin{tabular}{|l|p{7cm}|c|}
\hline
\textbf{Mesure} & \textbf{Implémentation} & \textbf{Statut} \\
\hline
Authentification JWT & Tokens signés HS256, expiration 1h & OK \\
\hline
Hachage mots de passe & bcryptjs, coût 12 & OK \\
\hline
Limitation de débit & 10 req/15min (login), 200 req/min (API) & OK \\
\hline
RBAC & Rôles admin/user, routes protégées & OK \\
\hline
PIN Kiosk & 4 chiffres, 3 tentatives max, verrouillage + alerte & OK \\
\hline
Accès externe & VPN Tailscale, aucun port public exposé & OK \\
\hline
HTTPS & Certificats Tailscale + terminaison TLS nginx & OK \\
\hline
Logs d'accès & Journalisation des actions sensibles en BDD & OK \\
\hline
\end{tabular}
\caption{Mesures de sécurité du système SENTYS}
\end{table}

\subsection{Défis techniques rencontrés}
% CORRECTION : les problèmes étaient en tête de section. Ils sont maintenant en fin de section 7.

\subsubsection{Compatibilité Linux vs Windows (casse des noms de fichiers)}
Le fichier \texttt{Loginpage.tsx} ne posait aucun problème sur Windows (NTFS insensible à la casse). Lors du déploiement sur Ubuntu (ext4, sensible à la casse), le build échouait. Solution : renommage via \texttt{git mv}.

\subsubsection{Verrou dpkg lors de l'installation}
Le processus \texttt{unattended-upgrades} tenait le verrou \texttt{/var/lib/dpkg/lock-frontend}. Solution : arrêt du service, désactivation des mises à jour automatiques.

\subsubsection{Permissions PostgreSQL}
L'erreur \texttt{permission denied for schema public} au premier démarrage. Solution : \texttt{GRANT ALL ON SCHEMA public} et \texttt{ALTER DEFAULT PRIVILEGES} en tant que superutilisateur.

\subsubsection{Module ffmpeg-static manquant en CI}
\texttt{ffmpeg-static} référencé dans le code mais absent du \texttt{package.json}. Absent après \texttt{npm ci --omit=dev}. Solution : ajout explicite aux dépendances de production.

\subsubsection{Transition ESP32-CAM vers Raspberry Pi Zero 2W}
Comme détaillé en section~5.1.1, les tests avec l'ESP32-CAM AI-Thinker ont révélé des limitations bloquantes : flux MJPEG instable, mode hors ligne irréalisable, brochage insuffisant pour le capteur PIR. La transition vers le Pi Zero 2W, plus onéreuse, a permis d'implémenter toutes les fonctionnalités de résilience requises.

\paragraph{Note sur la détection IA}
L'analyse vidéo par IA (Python + OpenCV/YOLO) est une évolution planifiée, volontairement exclue du périmètre du MVP dont l'objectif prioritaire est la stabilité du pipeline de capture, d'enregistrement et d'alerte.

\subsection{Outils de travail}

\begin{table}[H]
\centering
\small
\begin{tabular}{|l|l|p{6cm}|}
\hline
\textbf{Catégorie} & \textbf{Outil} & \textbf{Usage} \\
\hline
Versioning & Git + GitHub & Branches par feature, commits atomiques \\
\hline
CI/CD & GitHub Actions & Build + déploiement automatique sur push \\
\hline
Éditeur & VS Code & Développement React, Node.js \\
\hline
Base de données & PostgreSQL 16 & Stockage persistant \\
\hline
Gestion processus & PM2 & Démarrage auto + monitoring Node.js \\
\hline
Accès distant & Tailscale & VPN WireGuard sans configuration box \\
\hline
Documentation & \LaTeX{} & Rapport TFE complet \\
\hline
Assistance IA & Claude (Anthropic) & Mise en forme \LaTeX{}, reformulation et structuration \\
\hline
\end{tabular}
\caption{Environnement et outils de développement}
\end{table}

% ==========================================================
\section{Aspects Législatifs}
\label{sec:legal}
% SECTION NOUVELLE — CORRECTION : les obligations légales liées à l'installation
% de caméras en Belgique n'étaient pas traitées.
% ==========================================================

\subsection{Cadre légal belge sur la vidéosurveillance}

En Belgique, l'installation de caméras de surveillance est encadrée par la \textbf{loi du 21 mars 2007 réglant l'installation et l'utilisation de caméras de surveillance} (dite « Loi Caméras »), modifiée par la loi du 21 mars 2018 et complétée par le RGPD.

\subsubsection{Obligations déclaratives}

Depuis le 25 mai 2018 (entrée en vigueur du RGPD), les caméras installées dans un cadre privé et non accessible au public sont soumises aux obligations suivantes :

\begin{itemize}[label=\textbullet]
    \item \textbf{Déclaration à l'Autorité de Protection des Données (APD) :} conformément à l'article 30 du RGPD, le responsable du traitement doit tenir un registre des activités de traitement. Pour les installations résidentielles simples (exception domestique, Art. 2.2.c du RGPD), cette obligation est allégée si les caméras ne filment que la propriété privée du responsable.
    \item \textbf{Exception domestique :} le système SENTYS, configuré pour ne filmer que la propriété privée du client (entrée, garage, jardin), relève de cette exception. Aucune déclaration formelle à l'APD n'est requise pour cet usage strictement privé.
    \item \textbf{Limite de l'exception :} si le champ des caméras couvre la voie publique ou des propriétés voisines, l'exception cesse de s'appliquer et une déclaration à l'APD devient obligatoire.
\end{itemize}

\subsubsection{Signalétique obligatoire}

La loi impose l'affichage d'un pictogramme normalisé à l'entrée de toute zone surveillée, même dans un cadre privé, dès lors que des tiers (livreurs, visiteurs) peuvent y accéder. Ce pictogramme doit mentionner :

\begin{itemize}[label=\textbullet]
    \item L'existence de caméras de surveillance.
    \item L'identité du responsable de traitement.
    \item Le droit d'accès aux données.
\end{itemize}

Le client a été informé de cette obligation et s'est engagé à apposer la signalétique réglementaire à l'entrée de sa propriété.

\subsubsection{Durée de conservation}

La loi belge fixe une durée maximale de conservation des images à \textbf{30 jours}, sauf en cas d'incident en cours d'investigation. SENTYS implémente une purge automatique à 30 jours, en conformité avec cette exigence.

\subsubsection{Bilan de conformité}

\begin{table}[H]
\centering
\small
\renewcommand{\arraystretch}{1.3}
\begin{tabular}{|l|l|c|}
\hline
\textbf{Obligation} & \textbf{Mise en œuvre} & \textbf{Statut} \\
\hline
Champ limité à la propriété privée & Paramétrage des angles de vue lors de l'installation & \checkmark \\
\hline
Signalétique à l'entrée & Pictogramme fourni au client & \checkmark \\
\hline
Conservation $\leq$ 30 jours & Purge automatique implémentée & \checkmark \\
\hline
Aucun flux sur serveur tiers & Architecture 100\% locale (pas de cloud) & \checkmark \\
\hline
Droits d'accès/suppression & Interface d'administration (US-19) & \checkmark \\
\hline
\end{tabular}
\caption{Bilan de conformité légale du système SENTYS}
\end{table}

% ==========================================================
\section{Conclusion}
% ==========================================================

Le projet SENTYS constitue un système de surveillance résidentiel complet, entièrement local et résilient, répondant aux besoins identifiés dans le cahier des charges.

\subsection{Points forts}

\begin{itemize}[label=\textbullet]
    \item \textbf{Résilience totale :} coupure WiFi et panne électrique gérées indépendamment, sans perte de données, validé en conditions réelles chez le client.
    \item \textbf{Accès distant garanti :} VPN Tailscale fonctionnel depuis n'importe où, aucun port exposé publiquement.
    \item \textbf{CI/CD opérationnel :} chaque \texttt{git push} déclenche automatiquement le redéploiement.
    \item \textbf{PWA :} application installable sur smartphone, consultable hors ligne.
    \item \textbf{Conformité légale et RGPD :} privacy by design, conservation limitée à 30 jours, conformité à la loi belge sur les caméras, aucun cloud.
    \item \textbf{Coût maîtrisé :} moins de 250 euros de matériel, 0 € récurrent mensuel.
    \item \textbf{Stack moderne et éprouvé :} React/TypeScript, Node.js ESM, PostgreSQL, ffmpeg, PM2, nginx.
\end{itemize}

\subsection{Perspectives d'évolution}

\begin{itemize}[label=\textbullet]
    \item Intégration de la détection IA locale (Python + OpenCV/YOLO) : distinction personnes / animaux.
    \item Extension à d'autres capteurs (température, fumée, inondation).
    \item Interface multilingue (français / néerlandais / anglais).
    \item Tableau de bord analytique : statistiques d'activité par caméra et par période.
\end{itemize}

% ==========================================================
\section{Bibliographie}
% ==========================================================

\begin{enumerate}
    \item Espressif Systems. \textit{ESP32-CAM AI-Thinker datasheet}. Disponible sur : \url{https://www.espressif.com/}
    \item Tailscale Inc. \textit{Tailscale Documentation --- WireGuard-based VPN}. Disponible sur : \url{https://tailscale.com/kb/}
    \item GitHub. \textit{GitHub Actions --- Self-hosted runners}. Disponible sur : \url{https://docs.github.com/en/actions/hosting-your-own-runners}
    \item ffmpeg Project. \textit{ffmpeg Documentation}. Disponible sur : \url{https://ffmpeg.org/documentation.html}
    \item PM2. \textit{PM2 --- Advanced Node.js process manager}. Disponible sur : \url{https://pm2.keymetrics.io/}
    \item vite-plugin-pwa. \textit{Vite PWA Plugin Documentation}. Disponible sur : \url{https://vite-pwa-org.netlify.app/}
    \item Web Push Protocol. \textit{RFC 8030 --- Generic Event Delivery Using HTTP Push}. Disponible sur : \url{https://www.rfc-editor.org/rfc/rfc8030}
    \item Commission Européenne. \textit{Règlement Général sur la Protection des Données (RGPD) --- 2016/679}. Disponible sur : \url{https://eur-lex.europa.eu/}
    \item PostgreSQL Global Development Group. \textit{PostgreSQL 16 Documentation}. Disponible sur : \url{https://www.postgresql.org/docs/16/}
    \item Node.js Foundation. \textit{Node.js Documentation}. Disponible sur : \url{https://nodejs.org/docs/latest/api/}
    \item Autorité de Protection des Données (Belgique). \textit{Loi du 21 mars 2007 réglant l'installation et l'utilisation de caméras de surveillance}. Disponible sur : \url{https://www.autoriteprotectiondonnees.be/}
    \item OpenAI. \textit{GPT Models Documentation}. Disponible sur : \url{https://platform.openai.com/docs/}
    \item Anthropic. \textit{Claude AI Documentation}. Disponible sur : \url{https://docs.anthropic.com/}
    \item AlexxIT. \textit{go2rtc --- Stream application with multi-protocol support}. disponible sur : \url{https://github.com/AlexxIT/go2rtc}
\end{enumerate}

% ==========================================================
\appendix
\section{Annexes}
% ==========================================================

\subsection{Connexions matérielles du nœud caméra}

\begin{table}[H]
\centering
\begin{tabular}{|l|l|l|}
\hline
\textbf{Composant} & \textbf{Interface} & \textbf{Destination (RPi Zero 2W)} \\
\hline
Capteur InnoMaker OV5647 & Nappe FFC 15-pin & Port CSI \\
\hline
Shield UPS X306 & Pogo-Pins + I2C & GPIO 5V/GND + monitoring batterie \\
\hline
Batterie Li-ion & Connecteur PH2.0 & Entrée BAT (Shield X306) \\
\hline
\end{tabular}
\caption{Connexions matérielles du nœud caméra}
\end{table}

\subsection{Machine à états du nœud caméra}

\begin{table}[H]
\centering
\small
\begin{tabular}{|l|l|l|l|}
\hline
\textbf{État} & \textbf{Réseau} & \textbf{Comportement caméra} & \textbf{Déclencheur} \\
\hline
EN LIGNE & Serveur accessible & Stream RTSP vers serveur (MediaMTX) & Connexion serveur OK \\
\hline
HORS LIGNE & Serveur inaccessible & Enregistrement SD (clips 30s) & Perte connexion serveur \\
\hline
SYNCHRONISATION & Serveur de retour & Upload séquentiel des clips en attente & Reconnexion détectée \\
\hline
\end{tabular}
\caption{Machine à états du nœud caméra}
\end{table}

\subsection{Commandes de gestion du serveur}

\begin{lstlisting}[language=bash, caption=Commandes PM2 essentielles]
# Statut du serveur
pm2 status
# Logs en temps reel
pm2 logs SENTYS
# Redemarrer
pm2 restart SENTYS
\end{lstlisting}

\end{document}